\begin{figure}[tb]
  \centering
  %% reference
  \begin{subfigure}[t]{0.02\textwidth}
    % \vspace{.2cm}
    % \textbf{a}
    \raisebox{-1.25cm}{\rotatebox[origin=t]{90}{\textbf{Reference}}}
  \end{subfigure}
  \begin{subfigure}[t]{0.96\textwidth}
    \includegraphics[width=\linewidth,valign=t]{images/outlines_reference.png} \\
  \end{subfigure}\hfill
  %% Transformer
  \vspace{-.5cm}
  \begin{subfigure}[t]{0.02\textwidth}
    % \vspace{.2cm}
    % \textbf{b}
    \raisebox{-1.25cm}{\rotatebox[origin=t]{90}{\textbf{Transformer}}}
  \end{subfigure}
  \begin{subfigure}[t]{0.96\textwidth}
    \includegraphics[width=\linewidth,valign=t]{images/outlines_Transformer.png} \\
  \end{subfigure}\hfill
  %% single-step
  \vspace{-.5cm}
  \begin{subfigure}[t]{0.02\textwidth}
    % \vspace{.2cm}
    % \textbf{b}
    \raisebox{-1.25cm}{\rotatebox[origin=t]{90}{\textbf{Feedforward}}}
  \end{subfigure}
  \begin{subfigure}[t]{0.96\textwidth}
    \includegraphics[width=\linewidth,valign=t]{images/outlines_single_step.png} \\
  \end{subfigure}\hfill
  ~
    \caption{
    \textbf{(Prediction visualization)}
    A qualitative comparison of length-100 trajectories generated by the Trajectory Transformer and a feedforward Gaussian dynamics model from PETS, a state-of-the-art planning algorithm \cite{chua2018pets}.
    Both models were trained on trajectories collected by a single policy, for which a true trajectory is shown for reference.
    Compounding errors in the single-step model lead to physically implausible predictions, whereas the Transformer-generated trajectory is visually indistinguishable from those produced by the policy acting in the actual environment.
    The paths of the feet and head are traced through space for depiction of the movement between rendered frames.
    }
    \label{fig:humanoid}
    \vspace{-.3cm}
\end{figure}

\vspace{-.2cm}
\section{Experiments}
\label{sec:experiments}
\vspace{-.1cm}

Our experimental evaluation focuses on (1) the accuracy of the Trajectory Transformer as a long-horizon predictor compared to standard dynamics model parameterizations and (2) the utility of sequence modeling tools -- namely beam search -- as a control algorithm in the context of offline reinforcement learning, imitation learning, and goal-reaching.

\vspace{-.05cm}
\subsection{Model Analysis}
\vspace{-.05cm}

We begin by evaluating the Trajectory Transformer
as a long-horizon policy-conditioned predictive model.
The usual strategy for predicting trajectories given a policy is to rollout with a single-step model, with actions supplied by the policy.
Our protocol differs from the standard approach not only in that the model is not Markovian, but also in that it does not require access to a policy to make predictions -- the outputs of the policy are modeled alongside the states encountered by that policy.
Here, we focus only on the quality of the model's predictions; we use actions predicted by the model for an imitation learning method in the next subsection.

\begin{figure}
    \centering
    \includegraphics[width=1.0\linewidth]{images/error.pdf}
    \caption{
    \textbf{(Compounding model errors)}
    We compare the accuracy of the Trajectory Transformer (with uniform discretization) to that of the probabilistic feedforward model ensemble \citep{chua2018pets} over the course of a planning horizon in the humanoid environment, corresponding to the trajectories visualized in Figure~\ref{fig:humanoid}.
    The Trajectory Transformer has substantially better error compounding with respect to prediction horizon than the feedforward model.
    The discrete oracle is the maximum log likelihood attainable given the discretization size; see Appendix~\ref{app:oracle} for a discussion.
    }
    \label{fig:model_error}
    % \vspace{-.5cm}
\end{figure}

\paragraph{Trajectory predictions.}
Figure~\ref{fig:humanoid} depicts a visualization of predicted 100-timestep trajectories from our model after having trained on a dataset collected by a trained humanoid policy.
Though model-based methods have been applied to the humanoid task, prior works tend to keep the horizon intentionally short to prevent the accumulation of model errors \citep{janner2019mbpo,amos2020model}.
The reference model is the probabilistic ensemble implementation of PETS \citep{chua2018pets}; we tuned the number of models within the ensemble, the number of layers, and layer sizes, but were unable to produce a model that predicted accurate sequences for more than a few dozen steps.
In contrast, we see that the Trajectory Transformer's
long-horizon predictions are substantially more accurate, remaining visually indistinguishable from the ground-truth trajectories even after 100 predicted steps.
To our knowledge, no prior model-based RL algorithm has demonstrated predicted rollouts of such accuracy and length on tasks of comparable dimensionality.

\paragraph{Error accumulation.}
A quantitative account of the same finding is provided in Figure~\ref{fig:model_error}, in which we evaluate the model's accumulated error versus prediction horizon.
Standard predictive models tend to have excellent single-step errors but poor long-horizon accuracy, so instead of evaluating a test-set single-step likelihood, we sample 1000 trajectories from a fixed starting point to estimate the per-timestep state marginal predicted by each model.
We then report the likelihood of the states visited by the reference policy on a held-out set of trajectories under these predicted marginals.
To evaluate the likelihood under our discretized model, we treat each bin as a uniform distribution over its specified range; by construction, the model assigns zero probability outside of this range.

To better isolate the source of the Transformer's improved accuracy over standard single-step models, we also evaluate a Markovian variant of our same architecture.
This ablation has a truncated context window that prevents it from attending to more than one timestep in the past.
This model performs similarly to the trajectory Transformer on fully-observed environments, suggesting that architecture differences and increased expressivity from the autoregressive state discretization play a large role in the trajectory Transformer's long-horizon accuracy.
We construct a partially-observed version of the same humanoid environment, in which each dimension of every state is masked out with $50\%$ probability (Figure~\ref{fig:model_error} right), and find that, as expected, the long-horizon conditioning plays a larger role in the model's accuracy in this setting.

\paragraph{Attention patterns.}
We visualize the attention maps during model predictions in Figure~\ref{fig:attention}.
We find two primary attention patterns.
The first is a discovered Markovian strategy, in which a state prediction attends overwhelmingly to the previous transition.
The second is qualitatively striated, with the model attending to specific dimensions in multiple prior states for each state prediction.
Simultaneously, the action predictions attend to prior actions more than they do prior states.
The action dependencies contrast with the usual formulation of behavior cloning, in which actions are a function of only past states, but is reminiscent of the action filtering technique used in some planning algorithm to produce smoother action sequences \citep{nagabandi2019pddm}.

\begin{figure}
    \centering
    \includegraphics[width=0.35\linewidth]{images/markov.pdf}
    \hspace{0.15\linewidth}
    \includegraphics[width=0.35\linewidth]{images/striated.pdf}
    \begin{flushleft}
    \vspace{-5cm}
    ${\color{cred}\st}$ \hspace{6.75cm} ${\color{cred}\st}$ \\
    \vspace{.04cm}
    ${\color{cblue}\at}$ \hspace{6.75cm} ${\color{cblue}\at}$ \\
    \vspace{.95cm}
    {\color{cred}\Large\vdots} \hspace{6.75cm} \\
        \vspace{-0.95cm}\hspace{7.2cm}
        {\color{cred}\Large\vdots} \\
    \vspace{.1cm}
    {\color{cblue}\Large\vdots}  \hspace{6.75cm} \\
        \vspace{-0.95cm}\hspace{7.2cm}
        {\color{cblue}\Large\vdots} \\
    \vspace{.85cm}
    ${\color{cred}\bs_{t+5}}$  \hspace{6.4cm} ${\color{cred}\bs_{t+5}}$ \\
    \vspace{.02cm}
    ${\color{cblue}\ba_{t+5}}$ \hspace{6.4cm} ${\color{cblue}\ba_{t+5}}$ \\
    \vspace{.05cm}
    \vspace{.2cm}
    \end{flushleft}
    \vspace{.05cm}
    \caption{
    \textbf{(Attention patterns)}
    We observe two distinct types of attention masks during trajectory prediction.
    In the first, both {\color{cred}states} and {\color{cblue}actions} are dependent primarily on the immediately preceding transition, corresponding to a model that has learned the Markov property.
    The second strategy has a striated appearance, with state dimensions depending most strongly on the same dimension of multiple previous timesteps.
    Surprisingly, actions depend more on past actions than they do on past states, reminiscent of the action smoothing used in some trajectory optimization algorithms \citep{nagabandi2019pddm}.
    The above masks are produced by a first- and third-layer attention head during sequence prediction on the hopper benchmark; reward dimensions are omitted for this visualization.$^1$
    }
    \label{fig:attention}
\end{figure}

\subsection{Reinforcement Learning and Control}
\label{sec:exp_rl}

\paragraph{Offline reinforcement learning.}
We evaluate the Trajectory Transformer on a number of environments from the D4RL offline benchmark suite \citep{fu2020d4rl}, including the locomotion and AntMaze domains.
This evaluation is the most difficult of our control settings, as reward-maximizing behavior is the most qualitatively dissimilar from the types of behavior that are normally associated with unsupervised modeling -- namely, imitative behavior.
Results for the locomotion environments are shown in Table~\ref{table:d4rl}.
We compare against five other methods spanning other approaches to data-driven control: (1) behavior-regularized actor-critic (BRAC; \citealt{wu2019behavior}) and conservative $Q$-learning (CQL; \citealt{kumar2020conservative}) represent the current state-of-the-art in model-free offline RL; model-based offline planning (MBOP; \citealt{argenson2020model}) is the best-performing prior offline trajectory optimization technique; decision transformer (DT; \cite{chen2021decision}) is a concurrently-developed sequence-modeling approach that uses return-conditioning instead of planning; and behavior-cloning (BC) provides the performance of a pure imitative method.

The Trajectory Transformer performs on par with or better than all prior methods (Table~\ref{table:d4rl}).
The two discretization variants of the Trajectory Transformer, uniform and quantile, perform similarly on all environments except for HalfCheetah-Medium-Expert, where the large range of the velocities prevents the uniform discretization scheme from recovering the precise actuation required for enacting the expert policy.
As a result, the quantile discretization approach achieves a return of more than twice that of the uniform discretization.\blfootnote{
    $^1$More attention visualizations can be found at 
    \href{https://trajectory-transformer.github.io/attention}
    {\texttt{trajectory-transformer.github.io/attention}}
}

\paragraph{Combining with $Q$-functions.}
Though Monte Carlo value estimates are sufficient for many standard offline RL benchmarks, in sparse-reward and long-horizon settings they become too uninformative to guide the beam-search-based planning procedure.
In these problems, the value estimate from the Transformer can be replaced with a $Q$-function trained via dynamic programming.
We explore this combination by using the $Q$-function from the implicit $Q$-learning algorithm (IQL; \citealt{kostrikov2021implicit}) on the AntMaze navigation tasks \citep{fu2020d4rl}, for which there is only a sparse reward upon reaching the goal state.
These tasks evaluate temporal compositionality because they require stitching together multiple zero-reward trajectories in the dataset to reach a designated goal.

AntMaze results are provided in Table~\ref{table:antmaze}.
$Q$-guided Trajectory Transformer planning outperforms all prior methods on all maze sizes and dataset compositions.
In particular, it outperforms the IQL method from which we obtain the $Q$-function, underscoring that planning with a $Q$-function as a search heuristic can be less susceptible to errors in the $Q$-function than policy extraction.
However, because the $Q$-guided planning procedure still benefits from the temporal compositionality of both dynamic programming and planning, it outperforms return-conditioning approaches, such as the Decision Transformer, that suffer due to the lack of complete demonstrations in the AntMaze datasets.

\begin{table*}
\centering
\small
\begin{tabular*}{\textwidth}{@{\extracolsep{\fill}}llrrrrrrr}
\toprule
\multicolumn{1}{c}{\bf Dataset} & \multicolumn{1}{c}{\bf Environment} & \multicolumn{1}{c}{\bf BC} & \multicolumn{1}{c}{\bf MBOP} & \multicolumn{1}{c}{\bf BRAC} & \multicolumn{1}{c}{\bf CQL} & \multicolumn{1}{c}{\bf DT} & \multicolumn{1}{c}{\bf TT \scriptsize{\textcolor{cgrey}{(uniform)}}} & \multicolumn{1}{c}{\bf TT \scriptsize{\textcolor{cgrey}{(quantile)}}} \\ 
\midrule
Med-Expert & HalfCheetah & $59.9$ & $105.9$ & $41.9$ & $91.6$ & $86.8$ & $40.8$ \scriptsize{\raisebox{1pt}{$\pm 2.3$}} & $95.0$ \scriptsize{\raisebox{1pt}{$\pm 0.2$}} \\ 
Med-Expert & Hopper & $79.6$ & $55.1$ & $0.9$ & $105.4$ & $107.6$ & $106.0$ \scriptsize{\raisebox{1pt}{$\pm 0.28$}} & $110.0$ \scriptsize{\raisebox{1pt}{$\pm 2.7$}} \\ 
Med-Expert & Walker2d & $36.6$ & $70.2$ & $81.6$ & $108.8$ & $108.1$ & $91.0$ \scriptsize{\raisebox{1pt}{$\pm 2.8$}} & $101.9$ \scriptsize{\raisebox{1pt}{$\pm 6.8$}} \\ 
\midrule
Medium & HalfCheetah & $43.1$ & $44.6$ & $46.3$ & $44.0$ & $42.6$ & $44.0$ \scriptsize{\raisebox{1pt}{$\pm 0.31$}} & $46.9$ \scriptsize{\raisebox{1pt}{$\pm 0.4$}} \\ 
Medium & Hopper & $63.9$ & $48.8$ & $31.3$ & $58.5$ & $67.6$ & $67.4$ \scriptsize{\raisebox{1pt}{$\pm 2.9$}} & $61.1$ \scriptsize{\raisebox{1pt}{$\pm 3.6$}} \\ 
Medium & Walker2d & $77.3$ & $41.0$ & $81.1$ & $72.5$ & $74.0$ & $81.3$ \scriptsize{\raisebox{1pt}{$\pm 2.1$}} & $79.0$ \scriptsize{\raisebox{1pt}{$\pm 2.8$}} \\ 
\midrule
Med-Replay & HalfCheetah & $4.3$ & $42.3$ & $47.7$ & $45.5$ & $36.6$ & $44.1$ \scriptsize{\raisebox{1pt}{$\pm 0.9$}} & $41.9$ \scriptsize{\raisebox{1pt}{$\pm 2.5$}} \\ 
Med-Replay & Hopper & $27.6$ & $12.4$ & $0.6$ & $95.0$ & $82.7$ & $99.4$ \scriptsize{\raisebox{1pt}{$\pm 3.2$}} & $91.5$ \scriptsize{\raisebox{1pt}{$\pm 3.6$}} \\ 
Med-Replay & Walker2d & $36.9$ & $9.7$ & $0.9$ & $77.2$ & $66.6$ & $79.4$ \scriptsize{\raisebox{1pt}{$\pm 3.3$}} & $82.6$ \scriptsize{\raisebox{1pt}{$\pm 6.9$}} \\ 
\midrule
\multicolumn{2}{c}{\bf Average} & 47.7 & 47.8 & 36.9 & 77.6 & 74.7 & 72.6 \hspace{.6cm} & 78.9 \hspace{.6cm} \\ 
\bottomrule
\end{tabular*}
\vspace{.1cm}
\caption{
\textbf{(Offline reinforcement learning)} The Trajectory Transformer (TT) performs on par with or better than the best prior offline reinforcement learning algorithms on D4RL locomotion (\texttt{v2}) tasks.
Results for TT variants correspond to the mean and standard error over 15 random seeds (5 independently trained Transformers and 3 trajectories per Transformer).
We detail the sources of the performance for other methods in Appendix~\ref{app:baselines}.
}
\label{table:d4rl}
\end{table*}

\begin{figure}
\centering
\includegraphics[width=0.85\linewidth]{images/bar.pdf}
% \vspace{-0.1cm}
\caption{
\textbf{(Offline averages)}
A plot showing the average per-algorithm performance in Table~\ref{table:d4rl}, with bars colored according to a crude algorithm categorization.
In this plot, ``Trajectory Transformer'' refers to the quantile discreization variant.
}
\vspace{-.25cm}
\label{fig:offline}
\end{figure}


\begin{table*}[t]
\centering
\small
\begin{tabular*}{\textwidth}{@{\extracolsep{\fill}}llrrrrrr}
\toprule
\multicolumn{1}{c}{\bf Dataset} & \multicolumn{1}{c}{\bf Environment} &
    \multicolumn{1}{c}{\bf BC} &
    \multicolumn{1}{c}{\bf CQL} &
    \multicolumn{1}{c}{\bf IQL} &
    \multicolumn{1}{c}{\bf DT} &
    \multicolumn{1}{c}{\bf TT \scriptsize{\textcolor{cgrey}{$\bm{(+Q)}$}}} \\ 
\midrule
Umaze & AntMaze & 
    $54.6$ &
    $74.0$ &
    $87.5$  &
    $59.2$ &
    $100.0$ \scriptsize{\raisebox{1pt}{$\pm 0.0$}} \\ 
Medium-Play & AntMaze &
    $0.0$ &
    $61.2$ &
    $71.2$ &
    $0.0$ &
    $93.3$ \scriptsize{\raisebox{1pt}{$\pm 6.4$}} \\ 
Medium-Diverse & AntMaze &
    $0.0$ &
    $53.7$ &
    $70.0$ &
    $0.0$ &
    $100.0$ \scriptsize{\raisebox{1pt}{$\pm 0.0$}} \\
Large-Play & AntMaze &
    $0.0$ &
    $15.8$ &
    $39.6$ &
    $0.0$ &
    $66.7$ \scriptsize{\raisebox{1pt}{$\pm 12.2$}} \\ 
Large-Diverse & AntMaze &
    $0.0$ &
    $14.9$ &
    $47.5$ &
    $0.0$ &
    $60.0$ \scriptsize{\raisebox{1pt}{$\pm 12.7$}} \\
\midrule
\multicolumn{2}{c}{\bf Average} & 10.9 & 44.9 & 63.2 & 11.8 & 84.0 \hspace{.6cm} \\ 
\bottomrule
\end{tabular*}
\caption{
\textbf{(Combining with $Q$-functions)}
Performance on the sparse-reward AntMaze (\texttt{v0}) navigation task.
Using a $Q$-function as a search heuristic with the Trajectory Transformer ({TT~\scriptsize{\textcolor{cgrey}{$\bm{(+Q)}$}}}) outperforms policy extraction from the $Q$-function (IQL) and return-conditioning approaches like the Decision Transformer (DT).
We report means and standard error over 15 random seeds for {TT \scriptsize{\textcolor{cgrey}{$\bm{(+Q)}$}}}; baseline results are taken from \cite{kostrikov2021implicit}.
}
\vspace{-.1cm}
\label{table:antmaze}
\end{table*}



\paragraph{Imitation and goal-reaching.}
We additionally plan with the Trajectory Transformer using standard likelihood-maximizing beam search, as opposed to the return-maximizing version used for offline RL.
We find that after training the model on datasets collected by expert policies \citep{fu2020d4rl}, using beam search as a receding-horizon controller achieves an average normalized return of $104\%$ and $109\%$ in the Hopper and Walker2d environments, respectively, using the same evaluation protocol of 15 runs described as in the offline RL results.
While this result is perhaps unsurprising, as behavior cloning with standard feedforward architectures is already able to reproduce the behavior of the expert policies, it demonstrates that a decoding algorithm used for language modeling can be effectively repurposed for control.

Finally, we evaluate the goal-reaching variant of beam-search, which conditions on a future desired state alongside previously encountered states.
We use a continuous variant of the classic four rooms environment as a testbed \citep{sutton1999semimdps}.
Our training data consists of trajectories collected by a pretrained goal-reaching agent, with start and goal states sampled uniformly at random across the state space.
Figure~\ref{fig:goals} depicts routes taken by the the planner.
Anti-causal conditioning on a future state allows for beam search to be used as a goal-reaching method.
No reward shaping, or rewards of any sort, are required; the planning method relies entirely on goal relabeling.
An extension of this experiment to procedurally-generated maps is described in Appendix~\ref{app:minigrid}.

\begin{figure}[h]
    \centering
    \includegraphics[width=0.325\linewidth]{images/goals_red/0.png}
    \includegraphics[width=0.325\linewidth]{images/goals_red/1.png}
    \includegraphics[width=0.325\linewidth]{images/goals_red/2.png}
    \caption{
    \textbf{(Goal-reaching)}
    Trajectories collected by TTO with anti-causal goal-state conditioning in a continuous variant of the four rooms environment.
    Trajectories are visualized as curves passing through all encountered states, with color becoming more saturated as time progresses.
    Note that these curves depict real trajectories collected by the controller and not sampled sequences.
    The starting state is depicted by
    \protect{\includegraphics[height=.35cm]{images/goals_red/rolloutblack.pdf}}
    and the goal state by
    \protect{\includegraphics[height=.35cm]{images/goals_red/rolloutblue.pdf}}.
    Best viewed in color.
    }
    \vspace{-.2cm}
    \label{fig:goals}
\end{figure}

